% ---
% title: Two dice showing different numbers
% tags: conditional-probability
% ---

\begin{problem}
	Two fair dice are rolled. Given that they show different numbers, what is
	the conditional probability that (a) at least one of them is a six;
	(b) the sum of the two numbers is 7?
\end{problem}

\begin{solution}
	\idea Write the result as an ordered pair (first die, second die). All
	$36$ pairs are equally likely. The condition ``different numbers'' keeps
	only some of these pairs, and they stay equally likely. So we can count
	inside this reduced sample space. We also check the result with the
	formula $P(E\mid F)=P(EF)/P(F)$.
	\step{The condition} Let $F$ be the event ``the two dice show different
	numbers''. Its complement $F^c$ consists of the $6$ pairs with equal
	numbers: $(1,1),(2,2),\dots,(6,6)$. So
	\[ |F|=36-6=30, \qquad P(F)=\frac{30}{36}=\frac56. \]
	Knowing that $F$ happened, only these $30$ pairs are possible, and each of
	them is equally likely. This is the reduced sample space.
	\step{(a)} Let $E$ be the event ``at least one six''. We need the pairs in
	$F$ that contain a six, that is, the pairs in $EF$. Since the numbers are
	different, the pair cannot be $(6,6)$, so exactly one die shows a six and
	the other shows one of $1,\dots,5$:
	\begin{itemize}
		\item six on the first die: $(6,1),\dots,(6,5)$, that is, $5$ pairs;
		\item six on the second die: $(1,6),\dots,(5,6)$, that is, $5$ pairs.
	\end{itemize}
	These two groups have no common pair, so $|EF|=5+5=10$. Counting in the
	reduced sample space:
	\[ P(E\mid F)=\frac{|EF|}{|F|}=\frac{10}{30}=\frac13\approx0.3333. \]
	With the formula we get the same:
	\[ P(E\mid F)=\frac{P(EF)}{P(F)}=\frac{10/36}{30/36}=\frac{10}{30}=\frac13. \]
	For comparison, without the condition $E$ contains $11$ pairs (the $10$
	above and $(6,6)$), so $P(E)=11/36\approx0.3056$. The condition makes a six
	slightly more likely: it removes $5$ pairs without a six but only $1$ pair
	with a six.
	\step{(b)} Let $G$ be the event ``the sum is 7''. It has $6$ pairs:
	\[ (1,6),\ (2,5),\ (3,4),\ (4,3),\ (5,2),\ (6,1). \]
	In each of them the two numbers are different (two equal numbers give an
	even sum, and $7$ is odd). So all of $G$ lies inside $F$: $GF=G$ and
	$|GF|=6$. Counting in the reduced sample space:
	\[ P(G\mid F)=\frac{|GF|}{|F|}=\frac{6}{30}=\frac15=0.2. \]
	With the formula:
	\[ P(G\mid F)=\frac{P(GF)}{P(F)}=\frac{6/36}{30/36}=\frac{6}{30}=\frac15. \]
	Without the condition, $P(G)=6/36=1/6\approx0.1667$. The condition removes
	only pairs with a sum other than $7$, so the share of sum-$7$ pairs grows.
	\begin{answer}
		(a) $\dfrac13\approx0.3333$ \quad (b) $\dfrac15=0.2$
	\end{answer}
	\begin{remark}
		Counting in the reduced sample space works because all outcomes of $F$
		are equally likely. Then
		\[ \frac{P(EF)}{P(F)}=\frac{|EF|/36}{|F|/36}=\frac{|EF|}{|F|}, \]
		and the $36$ cancels. In (b) we have the special case $G\subseteq F$ of
		Exercise~\exref{special-cases}, so
		\[ P(G\mid F)=\frac{P(G)}{P(F)}=\frac{1/6}{5/6}=\frac15. \]

		\textbf{Common mistake:} dividing by $36$ instead of $30$. Once we know
		that the numbers are different, the $6$ pairs with equal numbers are
		no longer possible, so they must not be counted in the denominator.
	\end{remark}
\end{solution}
